 \section*{第 \ref{ch:g11:nucleus-radioactivity} 章——原子核与放射性}

\begin{solution}{exo:g11:nucleus-radioactivity:1}
${}^{16}_{8}\mathrm{O}$：$8$ 个质子，$8$ 个中子；${}^{27}_{13}\mathrm{Al}$：$13$ 个质子，
$14$ 个中子；${}^{60}_{27}\mathrm{Co}$：$27$ 个质子，$33$ 个中子；
${}^{235}_{92}\mathrm{U}$：$92$ 个质子，$143$ 个中子；${}^{3}_{1}\mathrm{H}$：$1$ 个质子，
$2$ 个中子。（每次均有 $N = A - Z$。）
\end{solution}

\begin{solution}{exo:g11:nucleus-radioactivity:2}
三个碳\omterm{def:g11:fundamental-interactions:ladder}{原子}（$Z = 6$）互为\omterm{def:g11:nucleus-radioactivity:notation}{同位素}。
${}^{14}_{6}\mathrm{C}$ 与 ${}^{14}_{7}\mathrm{N}$ 的 $A = 14$ 相同，但
$Z$ 不同，是不同元素（电子云不同），因此不是\omterm{def:g11:nucleus-radioactivity:notation}{同位素}。
\end{solution}

\begin{solution}{exo:g11:nucleus-radioactivity:3}
$A = 210 - 4 = 206$，$Z = 84 - 2 = 82$：
${}^{210}_{84}\mathrm{Po} \to {}^{206}_{82}\mathrm{Pb} +
{}^{4}_{2}\mathrm{He}$——子核为铅-206。
\end{solution}

\begin{solution}{exo:g11:nucleus-radioactivity:4}
${}^{60}_{27}\mathrm{Co} \to {}^{60}_{28}\mathrm{Ni} +
{}^{\;0}_{-1}\mathrm{e}$ 与
${}^{131}_{53}\mathrm{I} \to {}^{131}_{54}\mathrm{Xe} +
{}^{\;0}_{-1}\mathrm{e}$：两式中 $A$ 不变，$Z$ 增加 1。
\end{solution}

\begin{solution}{exo:g11:nucleus-radioactivity:5}
$16$ 天 $= 2\,T_{1/2}$：$800/4 = \qty{200}{MBq}$。$32$ 天
$= 4\,T_{1/2}$：$800/16 = \qty{50}{MBq}$。
\end{solution}

\begin{solution}{exo:g11:nucleus-radioactivity:6}
(a) $\alpha$（质量大，因而几乎不偏转）；(b) $\beta^-$；(c) $\gamma$
（中性，穿透力强）；(d) $\beta^+$（\omterm{rem:g11:nucleus-radioactivity:betaplus}{正电子}：与 (b) 一样轻，
电荷相反）。
\end{solution}

\begin{solution}{exo:g11:nucleus-radioactivity:7}
${}^{222}_{86}\mathrm{Rn} \to {}^{218}_{84}\mathrm{Po} +
{}^{4}_{2}\mathrm{He}$; \
${}^{218}_{84}\mathrm{Po} \to {}^{214}_{82}\mathrm{Pb} +
{}^{4}_{2}\mathrm{He}$; \
${}^{214}_{82}\mathrm{Pb} \to {}^{214}_{83}\mathrm{Bi} +
{}^{\;0}_{-1}\mathrm{e}$.
\end{solution}

\begin{solution}{exo:g11:nucleus-radioactivity:8}
氟-18 有 $9$ 个中子，而稳定的氟-19 有 $10$ 个中子：
中子数偏少，位于稳定带下方，因此发生 $\beta^+$ 衰变：
${}^{18}_{9}\mathrm{F} \to {}^{18}_{8}\mathrm{O} + {}^{0}_{+1}\mathrm{e}$.
\end{solution}

\begin{solution}{exo:g11:nucleus-radioactivity:9}
$24$ 小时 $= 4\,T_{1/2}$：$500/16 \approx \qty{31}{MBq}$。由于
$2^9 = 512$，九个\omterm{def:g11:nucleus-radioactivity:halflife}{半衰期}后 $500/512 < \qty{1}{MBq}$，经过
$9 \times 6 = \qty{54}{h}$.
\end{solution}

\begin{solution}{exo:g11:nucleus-radioactivity:10}
$\frac18 = \frac1{2^3}$：三个\omterm{def:g11:nucleus-radioactivity:halflife}{半衰期}，$3 \times \num{5700} =
\num{17100}$ 年。一块恐龙骨骼的年龄约为 $\num{17500}$ 个\omterm{def:g11:nucleus-radioactivity:halflife}{半衰期}：
剩下的份额为 $1/2^{17500}$——连一个碳-14 \omterm{def:g11:fundamental-interactions:ladder}{原子}也没有（一克碳
仅约含 $\num{5e22}$ 个\omterm{def:g11:fundamental-interactions:ladder}{原子}）。
\end{solution}

\begin{solution}{exo:g11:nucleus-radioactivity:11}
${}^{14}_{6}\mathrm{C}$：$8$ 个中子，比稳定碳中的 $6$--$7$ 个多，
位于稳定带上方，发生 $\beta^-$ 衰变。${}^{11}_{6}\mathrm{C}$：$5$ 个中子，
中子数偏少，发生 $\beta^+$ 衰变。${}^{238}_{92}\mathrm{U}$：$Z = 92 > 83$，过
重，发生 $\alpha$ 衰变。
\end{solution}

\begin{solution}{exo:g11:nucleus-radioactivity:12}
只有 $\alpha$ 衰变会改变 $A$：$238 - 206 = 32 = 4x$，故 $x = 8$ 次。它们
共使 $16$ 从 $Z$ 中减去：$92 - 16 + y = 82$，得 $y = 6$ 次负 β 衰变。
\end{solution}

\begin{solution}{exo:g11:nucleus-radioactivity:13}
$6.4/0.1 = 64 = 2^{6}$：六个\omterm{def:g11:nucleus-radioactivity:halflife}{半衰期}，$6 \times 5.3 = 31.8 \approx 32$
年，即在屏蔽条件下储存的时间。
\end{solution}

\begin{solution}{exo:g11:nucleus-radioactivity:14}
每天：$8000 \times 86400 \approx \num{6.9e8}$ 次衰变。经过 $80$ 年：
$\num{6.9e8} \times 365 \times 80 \approx \num{2e13}$。生命始终处于
这种本底辐射之中——一\omterm{def:g11:nucleus-radioactivity:activity}{贝可勒尔}即每秒一个\omterm{def:g11:fundamental-interactions:ladder}{原子}衰变，相对于人体的
$\sim 10^{27}$ 个\omterm{def:g11:fundamental-interactions:ladder}{原子}而言；\omterm{def:g10:orders-of-magnitude:unit}{单位}很小，听起来吓人的
\omterm{def:g11:nucleus-radioactivity:activity}{贝可勒尔}数字，可能只相当于几根香蕉的水平。
\end{solution}

\begin{solution}{exo:g11:nucleus-radioactivity:15}
每天约 $\num{25000} \times 86400 \approx \num{2.2e9}$ 次衰变。十年约为
一个\omterm{def:g11:nucleus-radioactivity:halflife}{半衰期}的 $10/432 \approx 2\%$：远少于一半已衰变，因此
迫使更换的是电子元件与灰尘，而非放射源。在外部，
$\alpha$ 粒子在腔室的空气和外壳中即被吸收；吸入的灰尘使放射源
停留在肺中，每个 $\alpha$ 粒子都把\omterm{def:g10:energy-conservation:energy}{能量}释放到活
体组织中——较小的\omterm{def:g11:nucleus-radioactivity:activity}{放射性活度}即可造成较大的\omterm{rem:g11:nucleus-radioactivity:dose}{剂量}（以\omterm{rem:g11:nucleus-radioactivity:dose}{希沃特}计）。
\end{solution}

\begin{solution}{pb:g11:nucleus-radioactivity:1}
\textbf{1.} ${}^{12}_{6}\mathrm{C}$：$6$ 个质子，$6$ 个中子；${}^{14}_{6}\mathrm{C}$：
$6$ 个质子，$8$ 个中子。化学只“看到”由 $Z = 6$ 决定的电子云：
二者化学行为相同。

\textbf{2.} 进食与呼吸不断更新生物体内的碳，
使其保持与大气相同的比例；死亡后摄入停止，衰变
便不受抵消地持续进行。

\textbf{3.} ${}^{14}_{6}\mathrm{C} \to {}^{14}_{7}\mathrm{N} +
{}^{\;0}_{-1}\mathrm{e}$（$\beta^-$）。这些 $\beta$ 粒子很“软”：几米
厚的裹尸亚麻布就能挡住它们，因此逃出木乃伊的很少——这正是为何
要测量从文物本身取下的一小份碳样品。

\textbf{4.} $6.8/13.6 = \frac12$：一个\omterm{def:g11:nucleus-radioactivity:halflife}{半衰期}。木乃伊的年龄
约为\num{5700}~年。

\textbf{5.} $13.2/13.6 \approx 0.97$：亚麻布基本属于现代——
至多几个世纪。这是一件赝品。

\textbf{6.} $13.6/2^{10} = 13.6/1024 \approx 0.013$ 为每分钟的衰变次数——
每克样品约每 $75$ 分钟才计得一次，淹没在天然本底中。十
个\omterm{def:g11:nucleus-radioactivity:halflife}{半衰期}（约\num{57000}~年）是实际的测年上限。

\textbf{7.} $92$ 个质子，$146$ 个中子。\omterm{def:g11:fundamental-interactions:charge}{库仑}斥力作用于
\omterm{def:g11:fundamental-interactions:ladder}{原子核}内所有质子对之间，而强核力只把
相邻\omterm{def:g11:fundamental-interactions:ladder}{核子}束缚在一起：重核只有依靠多余的中子才能存续，
即只增加“胶水”而不增加斥力。

\textbf{8.} ${}^{238}_{92}\mathrm{U} \to {}^{234}_{90}\mathrm{Th} +
{}^{4}_{2}\mathrm{He}$.

\textbf{9.} ${}^{234}_{90}\mathrm{Th} \to {}^{234}_{91}\mathrm{Pa} +
{}^{\;0}_{-1}\mathrm{e}$，然后 ${}^{234}_{91}\mathrm{Pa} \to
{}^{234}_{92}\mathrm{U} + {}^{\;0}_{-1}\mathrm{e}$：铀重新出现，即
铀-234。

\textbf{10.} $A$：$238 - 206 = 32$，故需 $8$ 次 $\alpha$ 衰变；$Z$：
$92 - 16 + y = 82$，故需 $6$ 次 $\beta^-$ 衰变。

\textbf{11.} $\num{4.5e9}$ 年为一个\omterm{def:g11:nucleus-radioactivity:halflife}{半衰期}：剩下 $\frac12$。再过五
十亿年 $\approx$ 再一个\omterm{def:g11:nucleus-radioactivity:halflife}{半衰期}：约剩 $\frac14$。

\textbf{12.} 它驱动火山活动、板块构造（大陆漂移、
地震）和地热：一颗由内部加热的行星，其表面地质
正源于此。

\textbf{13.} ${}^{241}_{95}\mathrm{Am} \to {}^{237}_{93}\mathrm{Np} +
{}^{4}_{2}\mathrm{He}$.

\textbf{14.} 每秒 $\num{3.3e4}$ 次衰变；
$\num{33000} \times 86400 \approx \num{2.9e9}$，即每天衰变次数。

\textbf{15.} 每个 $\alpha$ 粒子都会使所经过的空气电离；离子
携带微弱电流，穿过电离室两极之间。烟雾颗粒俘获
离子，电流减小，报警器响起。

\textbf{16.} $\alpha$ 粒子在电离室内耗尽全部\omterm{def:g10:energy-conservation:energy}{能量}，
没有一个逃出外壳：电离最强，漏出为零。一个等活度的 $\gamma$
源几乎不能使电离室内的空气电离，
而是照射整个房间。

\textbf{17.} 十年约占\omterm{def:g11:nucleus-radioactivity:halflife}{半衰期}的 $10/432 \approx 2\%$，
\omterm{def:g11:nucleus-radioactivity:activity}{放射性活度}几乎不变。薄弱环节是灰尘、昆虫和
电子器件——\omterm{def:g11:fundamental-interactions:ladder}{原子核}比这台装置更长寿。

\textbf{18.} $\num{9e9}$ 年 $= 2$ 个\omterm{def:g11:nucleus-radioactivity:halflife}{半衰期}：剩下 $\frac14$。
地球上的铀依然丰富，因此它至多形成于几个
\omterm{def:g11:nucleus-radioactivity:halflife}{半衰期}之前——地球是由新鲜的恒星灰烬凝聚而成，
即第一批恒星出现数十亿年之后。

\textbf{19.} 地质上已死亡：内部寒冷，
没有火山活动或板块构造来更新大气与岩石，核心也不再翻涌——因此
没有抵御太阳风的\omterm{def:g11:magnetic-fields:bfield}{磁场}屏障。

\textbf{20.} 你身上除氦以外的每个\omterm{def:g11:fundamental-interactions:ladder}{原子}都诞生于一颗垂死的恒星，
而那座“熔炉”的不稳定遗物仍在温暖你脚下
的大地——星尘因其自身的\omterm{def:g11:nucleus-radioactivity:halflife}{半衰期}而得以“延续”。
\end{solution}
